\documentclass[11pt]{article}
% Standalone preamble for web-only supplementary notes.
\usepackage[T1]{fontenc}
\usepackage{amsmath,amssymb,amsthm}
\usepackage{hyperref}
\usepackage{cleveref}
\DeclareMathOperator{\Map}{Map}
\DeclareMathOperator{\Fun}{Fun}
\newcommand{\iso}{\xrightarrow{\sim}}
\newtheorem{lemma}{Supplementary Lemma}
\newtheorem{corollary}[lemma]{Supplementary Corollary}
\crefname{lemma}{Supplementary Lemma}{Supplementary Lemmas}
\Crefname{lemma}{Supplementary Lemma}{Supplementary Lemmas}
\crefname{corollary}{Supplementary Corollary}{Supplementary Corollaries}
\Crefname{corollary}{Supplementary Corollary}{Supplementary Corollaries}

\begin{document}
\section*{Hom expressions and restriction}

The composition functors in the
\href{chapter-2-exercises.html}{elementary Yoneda exercise} are built by
passing between functors into a hom category and arrow families with fixed
endpoints. This note explains the normalization involved in that passage.

Assume the Chapter 1 structures, including pullbacks and functor categories,
and the walking morphism. Fix an absolute category $C$ and absolute objects $x,y$. Put
$H=\operatorname{Hom}_C(x,y)$, the pullback of endpoint evaluation
$\operatorname{Ar}(C)\to C\times C$ along $(x,y)\colon\star\to C\times C$.
A framed expression over an absolute category $\Gamma$ is a functor
$A\colon\Gamma\to\operatorname{Ar}(C)$ with specified comparisons
$p\colon e_0A=x!_\Gamma$ and $q\colon e_1A=y!_\Gamma$.
An endpoint-preserving comparison is an identification of arrow families
compatible with both frames, as in
\href{supplementary-arrows-with-specified-endpoints.html}{Arrows with specified endpoints}.
We write $\simeq$ between expressions for such comparison data.
No category of all these expressions is required.

\begin{lemma}
\label[lemma]{lem:Supplementary_Hom_Normalization}
There are operations $\operatorname{intro}(f)\colon\Gamma\to H$ for
framed expressions $f$, and $\operatorname{expr}(h)$ for functors
$h\colon\Gamma\to H$, with comparisons
\[
 \operatorname{intro}(\operatorname{expr}(h))=h,
 \qquad \operatorname{expr}(\operatorname{intro}(f))\simeq f.
\]
An identification of functors induces an endpoint-preserving comparison
of their expressions; a comparison of expressions induces an identification
of their introduction functors.
\end{lemma}
\begin{proof}[Proof idea]
Introduction lifts the endpoint cone with terminal leg $!_\Gamma$.
For extraction, compose $h$ with the arrow projection of $H$, and restrict
the two universal endpoint frames. Their base leg is initially the
composite of $h$ with the projection $H\to\star$. Normalize it to
$!_\Gamma$ using the terminal comparison.

Compare the original restricted cone and this normalized cone by the
identity on the arrow leg and the terminal comparison on the base leg.
The two projected matching equations verify a full cone comparison.
Pullback reconstruction gives the first formula. Pullback computation,
together with the inverse normalization comparison, gives the second;
terminal uniqueness aligns its base comparison and retains both endpoint equations.

Naturality of the original frames and of terminal normalization gives the
assertion for identifications of functors. In the reverse direction, pair
the two endpoint equations to form a cone comparison, then lift it through
the pullback property.
\end{proof}

\begin{lemma}
\label[lemma]{lem:Supplementary_Hom_Restriction}
For $r\colon\Delta\to\Gamma$, let $r^*f$ be restriction of a framed
expression, followed by normalization of its constant endpoints to
$x!_\Delta,y!_\Delta$. Then
\[
 \operatorname{intro}(f)r=\operatorname{intro}(r^*f),
 \qquad \operatorname{expr}(hr)\simeq r^*\operatorname{expr}(h).
\]
\end{lemma}
\begin{proof}[Proof idea]
The restricted endpoint cone has terminal leg $!_\Gamma r$; the cone
encoded from $r^*f$ has terminal leg $!_\Delta$. Compare them by the
identity on the arrow leg and the terminal identification on the base.
Each matching coordinate includes the associator from restriction and
the comparison between the two constant endpoints. Checking these two
coordinates gives the whole cone comparison.
Lift it and use compatibility of pullback factorization with restriction.
This gives the first formula. Apply extraction and the two reconstruction
comparisons of \Cref{lem:Supplementary_Hom_Normalization} to obtain the second.
\end{proof}

\begin{lemma}
\label[lemma]{lem:Supplementary_Hom_Composition_Computation}
Assume in addition the interval endpoint data and the Segal axiom.
For $e\colon\star\to\operatorname{Hom}_C(x,y)$, let
$f_\Gamma=\operatorname{expr}(e!_\Gamma)$.
For every absolute object $z$ of $C$, there is a postcomposition functor
$e_*\colon\operatorname{Hom}_C(z,x)\to\operatorname{Hom}_C(z,y)$
with an endpoint-preserving computation comparison
\[
 \operatorname{expr}(e_*h)\simeq f_\Gamma\circ\operatorname{expr}(h)
\]
for every $h\colon\Gamma\to\operatorname{Hom}_C(z,x)$.
There is a corresponding precomposition functor
$e^*\colon\operatorname{Hom}_C(y,z)\to\operatorname{Hom}_C(x,z)$,
with $\operatorname{expr}(e^*k)\simeq\operatorname{expr}(k)\circ f_\Gamma$
for every $k\colon\Gamma\to\operatorname{Hom}_C(y,z)$.
\end{lemma}
\begin{proof}[Proof idea]
For postcomposition, use the universal framed expression on
$\operatorname{Hom}_C(z,x)$, compose it with the constant family of $e$,
and apply introduction. Restrict this construction by $h$.
\Cref{lem:Supplementary_Hom_Restriction} moves restriction through
introduction and extraction. The
\href{supplementary-substitution-in-morphism-expressions.html}{substitution comparison
for composition} identifies the restricted composite with the composite
of the restricted inputs. Normalize the constant input and use the
introduction computation law. The result preserves both endpoints.
Precomposition uses the same construction with the other input fixed.
\end{proof}
\end{document}
